\section{Introduction}
\label{sec:intro}

A clinician deciding whether to order a lumbar puncture, a tau PET scan, or an anti-amyloid therapy
cannot act on a probability alone. This is why Alzheimer's-disease (AD) biomarker models keep
getting more accurate and keep not being deployed. The field's standard answer is \emph{post-hoc}
attribution---SHAP over a gradient-boosted tree or a tabular foundation model, exactly what
\citet{ding2025l2ctabpfn} provide alongside their state-of-the-art progression predictor. Post-hoc
attributions explain the model; they do not explain a line of reasoning, and they are known to be
unstable.

Chain-of-thought (CoT) prompting appears to offer something categorically different: an
\emph{ante-hoc}, natural-language rationale produced by the same forward pass that produces the
answer~\citep{wei2022cot}. If that rationale is causally connected to the answer, it is a genuinely
new kind of clinical interpretability---a rationale a clinician can read, check against the
patient's chart, and disagree with. If it is not, then deploying CoT in a clinical biomarker
pipeline manufactures unwarranted trust, which is a strictly worse outcome than an honest black box.
Our main question is which of these two is true.

\para{The gap.} Nobody has measured it. The AD tabular-LLM literature asserts interpretability
without testing it: \citet{kearney2026tapgpt}, the closest prior work, state explicitly that they
``did not enable extended thinking or reasoning modes,'' and assess interpretability qualitatively.
The CoT-faithfulness literature has built rigorous causal tests---early answering, mistake
injection, paraphrasing, filler tokens~\citep{lanham2023faithfulness}---but has applied them only to
generic multiple-choice NLP benchmarks, where an unfaithful rationale costs nothing. Worse, every
one of those tests assumes a discrete answer set: ``did the answer flip?'' is undefined for brain
age in years, so CoT for clinical \emph{regression} is entirely unmeasured.

\para{What we do.} We run three AD prediction tasks built from open cohorts---neuropathological tau
staging from real quantified A$\beta$40/A$\beta$42/tau/pTau assays (\taskone), brain-age regression
on cognitively normal scans (\tasktwo), and longitudinal dementia progression under the
longitudinal-to-cross-sectional transform of \citet{ding2025l2ctabpfn} (\taskthree)---and we score
every arm on accuracy, causal faithfulness, and robustness under the same grouped cross-validation,
the same feature list, and the same perturbations. Critically, CoT is compared not only against
tabular black boxes but against two matched LLM controls: a \direct condition (same prompt, no
reasoning) and a compute-matched \filler condition (the assistant turn is prefilled with a run of
``\texttt{.}'' tokens as long as the measured median chain for that task). Without both controls,
``CoT helps'' is untestable.

\para{What we find.} CoT is worse, not better. It trails the best tabular baseline by
$\Delta$AUC $-0.18$ on \taskone and $-0.16$ on \taskthree, and on \tasktwo its 11.43\,y MAE is worse
than the 9.07\,y of simply predicting the cohort mean age. On \taskone it is significantly
\emph{worse} than semantically empty filler tokens of identical length ($\Delta$AUC $-0.090$,
$q=0.048$). And the chains are largely inert: showing the model none of its own reasoning and
forcing an answer reproduces the full-chain answer 82.9\% of the time (\aoc $=0.066$), while a
10$\times$ corruption of a biomarker value inside the chain changes the answer 3.6\% of the time.
A blinded rubric with a validated shuffled-explanation control rates those same chains 3.67/4. The
chains are persuasive in proportion to how little they do.

Three results cut against a blanket dismissal, and they are the reason this is more than a null
result. Faithfulness is 5.7$\times$ higher on the continuous target than on classification
(\aoc $0.378$ vs.\ $0.066$). Our new Cited-Feature Intervention Test shows that the specific
biomarkers a chain names shift its prediction 1.4--2.1$\times$ more than the ones it ignores. And
the effect reverses with scale: at 1.5B parameters CoT \emph{improves} accuracy
($\Delta$AUC $+0.082$, $p=0.012$) and its chains are 3.7$\times$ more load-bearing than the 7B
model's---CoT is faithful exactly where the model needs it to compute, and decorative exactly where
it does not, which is backwards from the deployment incentive.

In summary, our main contributions are:
\begin{itemize}[leftmargin=*,itemsep=0pt,topsep=0pt]
    \item We provide the first measurement of CoT \emph{causal} faithfulness on clinical biomarker
    tasks, transplanting the intervention battery of \citet{lanham2023faithfulness} onto real
    quantified amyloid/tau neuropathology and two AD cohort tasks.
    \item We redefine that battery for a continuous target, replacing answer-flip rate with a
    normalised answer shift $|\Delta\hat{y}|/\mathrm{SD}(\hat{y})$, and use it to show that CoT
    faithfulness is systematically higher for regression than for classification.
    \item We propose two automatable interpretability metrics: \cfit, which perturbs the features a
    chain cites against the ones it ignores, and \esa, which rank-correlates citation frequency
    against independently measured permutation importance.
    \item We conduct a controlled comparison against black-box baselines \emph{and} two matched
    LLM controls, and document a dissociation---3.67/4 rated plausibility with 3.6\% causal
    sensitivity---that we argue should disqualify CoT as a clinical interpretability layer unless
    faithfulness is reported alongside accuracy.
\end{itemize}

\secref{sec:related} positions this against prior work, \secref{sec:method} describes the tasks,
arms, and metrics, \secref{sec:results} reports accuracy, faithfulness, robustness, and the clinical
brain-age endpoint, and \secref{sec:discussion} discusses what survives and what does not.
